\section{Check reconstruction independently of the encoder}
An encoder and decoder can share a component-order mistake and still round-trip their
own bytes. Compare decoded RGBA values with a separately constructed pixel array, and
open encoded results with an independent image decoder. For the sprite, derive that
array from the original eight row bytes rather than from the candidate's output.
For the eight-pixel trace, compare every alpha value as well as RGB.

Exercise each command at its boundaries. Include initial transparent-black INDEX,
wraparound DIFF values, LUMA extremes, collisions that replace an earlier cache entry,
runs of one and sixty-two, and runs that would exceed the declared image. Truncate
every multi-byte command at each missing-byte position. Give the last valid pixel a
wrong marker, an extra payload byte and a trailing byte as separate malformed cases.
Supply short destination capacities and counts whose high bit is set.

Run a failure and a different valid image in the same session. The valid image must
not inherit a previous color, a pending run or a cache entry from the failed operation.
For file routines, add forced short reads and writes, nonstandard descriptor numbers,
EOF at exact capacity and errors after a partial transfer. The QOI format is different
from PPM, but the descriptor ownership and accepted-byte-count reasoning still applies.

A permissive external decoder is evidence for successful interoperability, not the
exercise's malformed-input oracle. It may accept trailing bytes or a damaged marker.
Use the published strict contract when checking rejection and unchanged output. Keep
those decisions visible rather than treating another library's behavior as the format
specification.
